\section{Extended RQ2 Loop Trace Metrics}\label{app:rq2-extended}

\paragraph{Loop trace metric definitions for Table~\ref{tab:rq2-grid}.}
\emph{Planning fidelity.}
\textbf{Edge~F1} = F1 between the plan DAG edge set in the recorded loop trace and the human development DAG edge set.
\textbf{Layer~$\rho$} = Spearman correlation between topological layer indices in the recorded loop trace and those in the human development DAG. \textbf{CPR} = critical path length of $G_A$ divided by $|V_A|$ (direct value, unnormalized to gold). \textbf{WR} = mean ready frontier width of $G_A$ divided by that of $G_H$.
\emph{Implementation} metrics are computed on units resolved by the evaluated loop. \textbf{PatchLen} = patch lines per matched unit divided by gold patch lines. \textbf{Jacc} = token level Jaccard distance between candidate patches in the recorded loop trace and gold patch.
\emph{Testing} metrics summarize verification work authored during the recorded loop trace.
\textbf{\#T} = number of tests authored in the recorded loop trace per task.
\textbf{F2P} = fraction of tests authored in the recorded loop trace that fail on the initial state and pass on the gold solution.
\paragraph{Auxiliary planning metrics.}
Table~\ref{tab:rq2-grid} presents a compact metric set for planning, implementation, and testing. The auxiliary metrics here either duplicate headline columns or depend more directly on task graph shape than on loop implementation. Their appendix placement keeps the main table focused on loop trace metrics that remain comparable across loop implementations.

\paragraph{TC-F1 (planning).}
Transitive closure F1 (Appendix~\ref{app:tc-f1}) is strongly correlated with Edge F1 (Appendix~\ref{app:edge-f1}) in our runs ($\rho{>}0.9$ across the six evaluated loops), so it provides secondary evidence beyond the Edge F1 column in the main table. Its absolute level estimates whether the recorded plan captures \emph{indirect} prerequisite structure relevant to routing and obligation retention.

\paragraph{Branch Switching Rate (planning).}
BSR (Appendix~\ref{app:branch-switch}) is graph conditional. Higher BSR is consistent with frontier exploration when $G_H$ has many simultaneously ready branches, but on long sequential chains it can instead record unnecessary alternation across branches. We therefore report it as an extended loop trace metric outside the main table ranking columns.

\paragraph{Test-Ratio (testing).}
The main table reports \#T as an absolute test count. Test-Ratio equals $\#T/\#T_{\text{native}}$ with $\#T_{\text{native}}=44$, making it directly derivable from the \#T column and secondary to the reported verification effort metric.

\paragraph{JSD (testing).}
JSD measures distributional divergence between tests authored during the recorded loop trace and native test placement. Appendix placement reflects its auxiliary role relative to the main \#T and F2P columns, which more directly measure verification effort and fail to pass yield during loop engineering evaluation. Unit position histograms separate early unit concentration from late unit concentration, capturing where verification effort is routed along the horizon.

\paragraph{Source of additional patch lines.}
The main text reports that PatchLen exceeds the gold reference while Jacc to gold stays at a comparable moderate level. The recorded diffs explain the two metrics jointly. Candidate patches often impose a uniform style, add extensive comments, and rewrite existing repository blocks when repository state in the loop is incomplete. This pattern increases patch length while token level overlap remains close to the human reference, suggesting that loop state gaps can appear as implementation surplus rather than disjoint implementations.
